\documentclass[handout]{beamer}
\mode<presentation>
\usepackage{xcolor}
\usepackage[toc]{multitoc}
\usepackage{color}
\usepackage{graphicx}
\graphicspath{{../buku_ajar/buku_ajar_logika_matematika/figures/}}
\usepackage{tikz}
\usetikzlibrary{shapes,arrows,positioning,calc}
\usepackage{listings}
\usepackage{lmodern}
\usetheme[secheader]{Boadilla}
\usecolortheme{whale}
\setbeamertemplate{caption}[numbered]
\setbeamertemplate{bibliography item}[text]
\setbeamertemplate{navigation symbols}{}

\theoremstyle{definition}
\newtheorem{definisi}[theorem]{Definisi}
\theoremstyle{example}
\newtheorem{contoh}[theorem]{Contoh}

\renewcommand*{\multicolumntoc}{2}

\setbeamertemplate{footline}
{
  \leavevmode%
  \hbox{%
  \begin{beamercolorbox}[wd=.333333\paperwidth,ht=2.25ex,dp=1ex,center]{author in head/foot}%
    \usebeamerfont{author in head/foot}\insertshortauthor%
  \end{beamercolorbox}%
  \begin{beamercolorbox}[wd=.433333\paperwidth,ht=2.25ex,dp=1ex,center]{title in head/foot}%
    \usebeamerfont{title in head/foot}\insertshorttitle
  \end{beamercolorbox}%
  \begin{beamercolorbox}[wd=.233333\paperwidth,ht=2.25ex,dp=1ex,right]{date in head/foot}%
    \usebeamerfont{date in head/foot}\insertshortdate{}\hspace*{2em} \insertframenumber{} / \inserttotalframenumber\hspace*{2ex} 
  \end{beamercolorbox}}%
  \vskip0pt%
}

\subtitle[Logika Matematika]{Logika Matematika}
\title[Proposisi dan Kata Hubung Logika]{Konsep Dasar Proposisi dan Kata Hubung Logika}
\author{Cecep Suwanda, S.Si., M.Kom.}
\date{Maret 2025}
\institute{Universitas Bale Bandung}

\begin{document}

\maketitle

\section*{Daftar Isi}
\begin{frame}{Daftar Isi}
\frametitle{Daftar Isi}
\tableofcontents
\end{frame}

% ============================================================
% SECTION 3-1: Pengantar Proposisi
% ============================================================
\section{Pengantar Proposisi dan Kalimat Deklaratif}

\begin{frame}
\frametitle{Dari Bahasa ke Logika}
\begin{itemize}
  \item Mesin hanya beroperasi pada nilai \textbf{benar (True)} dan \textbf{salah (False)}.
  \item Bahasa alamiah perlu disaring menjadi unit pernyataan dengan nilai kebenaran tegas.
  \item Unit elementer tersebut disebut \textbf{proposisi}.
\end{itemize}
\end{frame}

\begin{frame}
\frametitle{Definisi Proposisi}
\begin{definisi}[Proposisi / Kalimat Tertutup]
\textbf{Proposisi} adalah \textbf{kalimat deklaratif} yang \textbf{hanya memiliki satu nilai kebenaran}: benar (T) atau salah (F), dan tidak mungkin benar sekaligus salah pada konteks yang sama.\\
Ketegasan ini sejalan dengan \textit{Law of Excluded Middle}.
\end{definisi}
\end{frame}

\begin{frame}
\frametitle{Contoh Proposisi}
\begin{contoh}
\begin{enumerate}
  \item ``Menara Eiffel terletak di pusat ibu kota Jepang.'' $\rightarrow$ Proposisi (salah secara geografis).
  \item ``Tiga ditambah dua sama dengan lima.'' $\rightarrow$ Proposisi benar.
  \item ``Kucing bernapas menggunakan paru-paru.'' $\rightarrow$ Proposisi (divalidasi biologi).
\end{enumerate}
\end{contoh}
\end{frame}

\begin{frame}
\frametitle{Kategori Non-Proposisi}
\begin{contoh}[Bukan Proposisi]
\begin{itemize}
  \item \textbf{Kalimat tanya:} ``Berapakah suhu di Mars?'' --- meminta informasi, bukan klaim.
  \item \textbf{Kalimat perintah:} ``Simpan berkas tersebut.'' --- tidak memiliki nilai kebenaran.
  \item \textbf{Kalimat seru:} ``Alangkah indahnya patung itu!'' --- ungkapan subjektif.
\end{itemize}
\end{contoh}
\end{frame}

\begin{frame}
\frametitle{Kalimat Terbuka}
\begin{definisi}[Kalimat Terbuka]
Kalimat yang \textbf{mengandung variabel bebas}, sehingga nilai kebenarannya belum dapat ditentukan hingga variabel disubstitusi.
\end{definisi}
\begin{contoh}
\begin{itemize}
  \item ``$x + 7 = 12$'' --- jika $x = 5$ benar; jika $x = 10$ salah.
  \item ``Dia adalah mahasiswa teladan.'' --- kata ``dia'' belum merujuk individu tertentu.
\end{itemize}
\end{contoh}
\end{frame}

% ============================================================
% SECTION 3-2: Epistemologi Nilai Kebenaran
% ============================================================
\section{Epistemologi Nilai Kebenaran}

\begin{frame}
\frametitle{Sumber Nilai Kebenaran}
\begin{itemize}
  \item Notasi: $\tau(p) = \mathbf{T}$ berarti ``proposisi $p$ bernilai benar''.
  \item Dua jalur epistemologis:
  \begin{itemize}
    \item \textbf{Observasi empiris} --- pengamatan fakta di dunia nyata.
    \item \textbf{Derivasi deduktif} --- dari definisi, aksioma, dan aturan inferensi.
  \end{itemize}
\end{itemize}
\end{frame}

\begin{frame}
\frametitle{Dasar Empiris}
\begin{contoh}
$q$: ``Air pada tekanan atmosfer di Himalaya tidak mendidih pada $80^\circ$C.''\\
Untuk menetapkan $\tau(q)$: pengamatan dan pengukuran (percobaan pemanasan air), bukan sekadar rumus tanpa konteks.
\end{contoh}
\end{frame}

\begin{frame}
\frametitle{Dasar Derivasi (Rasionalisme Matematis)}
\begin{contoh}
$r$: ``Penyelesaian $3x - 1 = 5$ adalah $x = 2$.''\\
Bukti deduktif: $3x - 1 = 5 \Rightarrow 3x = 6 \Rightarrow x = 2$.\\
Tanpa observasi lapangan, kesimpulan diperoleh dari manipulasi aljabar.
\end{contoh}
\end{frame}

% ============================================================
% SECTION 3-3: Operator Fundamental
% ============================================================
\section{Operator Fundamental: Negasi, Konjungsi, Disjungsi}

\begin{frame}
\frametitle{Negasi ($\neg$)}
\begin{definisi}[Negasi]
Operator uner: $\neg p$ bernilai kebalikan dari $p$.\\
Jika $p$ = T maka $\neg p$ = F; jika $p$ = F maka $\neg p$ = T.
\end{definisi}
\begin{block}{Tabel Kebenaran}
\begin{tabular}{|c|c|}
\hline
$p$ & $\neg p$ \\ \hline
T & F \\ \hline
F & T \\ \hline
\end{tabular}
\end{block}
\end{frame}

\begin{frame}
\frametitle{Konjungsi ($\wedge$) --- AND}
\begin{definisi}[Konjungsi]
$p \wedge q$ bernilai benar \textbf{hanya jika} $p$ dan $q$ keduanya benar.
\end{definisi}
\begin{block}{Tabel Kebenaran}
\begin{tabular}{|c|c|c|}
\hline
$p$ & $q$ & $p \wedge q$ \\ \hline
T & T & T \\ \hline
T & F & F \\ \hline
F & T & F \\ \hline
F & F & F \\ \hline
\end{tabular}
\end{block}
\end{frame}

\begin{frame}
\frametitle{Disjungsi ($\vee$) --- OR}
\begin{definisi}[Disjungsi Inklusif]
$p \vee q$ bernilai benar jika \textbf{minimal salah satu} $p$ atau $q$ benar.\\
Bernilai salah hanya jika keduanya salah.
\end{definisi}
\begin{block}{Tabel Kebenaran}
\begin{tabular}{|c|c|c|}
\hline
$p$ & $q$ & $p \vee q$ \\ \hline
T & T & T \\ \hline
T & F & T \\ \hline
F & T & T \\ \hline
F & F & F \\ \hline
\end{tabular}
\end{block}
\end{frame}

% ============================================================
% SECTION 3-4: Implikasi dan Biimplikasi
% ============================================================
\section{Implikasi dan Biimplikasi}

\begin{frame}
\frametitle{Implikasi ($\rightarrow$) --- Jika ... Maka ...}
\begin{itemize}
  \item $p \rightarrow q$: $p$ = anteseden (hipotesis), $q$ = konsekuen.
  \item Berpadanan dengan \texttt{IF ... THEN ...} dalam pemrograman.
\end{itemize}
\begin{definisi}
$p \rightarrow q$ bernilai salah \textbf{hanya} ketika $p$ = T dan $q$ = F.\\
Untuk kasus lain (termasuk $p$ = F), bernilai benar.
\end{definisi}
\end{frame}

\begin{frame}
\frametitle{Vacuous Truth}
\begin{block}{Kasus Implikasi}
\begin{itemize}
  \item \textbf{$T \rightarrow T$:} Janji dipenuhi $\Rightarrow$ benar.
  \item \textbf{$T \rightarrow F$:} Janji dilanggar $\Rightarrow$ salah.
  \item \textbf{$F \rightarrow T$, $F \rightarrow F$:} Anteseden tidak terpenuhi $\Rightarrow$ \textit{vacuous truth} (benar).
\end{itemize}
\end{block}
\end{frame}

\begin{frame}
\frametitle{Biimplikasi ($\leftrightarrow$) --- Jika dan Hanya Jika}
\begin{definisi}[Biimplikasi]
$p \leftrightarrow q$ bernilai benar jika $p$ dan $q$ sama nilainya (keduanya T atau keduanya F). Bernilai salah jika berbeda.
\end{definisi}
\begin{block}{Tabel Kebenaran}
\begin{tabular}{|c|c|c|}
\hline
$p$ & $q$ & $p \leftrightarrow q$ \\ \hline
T & T & T \\ \hline
T & F & F \\ \hline
F & T & F \\ \hline
F & F & T \\ \hline
\end{tabular}
\end{block}
\end{frame}

% ============================================================
% SECTION 3-5: Varian Implikasi dan Negasi Majemuk
% ============================================================
\section{Varian Implikasi dan Negasi Majemuk}

\begin{frame}
\frametitle{Konvers, Invers, Kontraposisi}
\begin{block}{Dari $p \rightarrow q$:}
\begin{itemize}
  \item \textbf{Konvers:} $q \rightarrow p$ --- \textbf{tidak ekuivalen}
  \item \textbf{Invers:} $\neg p \rightarrow \neg q$ --- \textbf{tidak ekuivalen}
  \item \textbf{Kontraposisi:} $\neg q \rightarrow \neg p$ --- \textbf{ekuivalen} dengan implikasi asli
\end{itemize}
\end{block}
\[
p \rightarrow q \equiv \neg q \rightarrow \neg p
\]
\end{frame}

\begin{frame}
\frametitle{Hukum De Morgan}
\begin{block}{Negasi Konjungsi dan Disjungsi}
\begin{align*}
\neg(p \wedge q) &\equiv \neg p \vee \neg q \\
\neg(p \vee q) &\equiv \neg p \wedge \neg q
\end{align*}
\end{block}
\end{frame}

\begin{frame}
\frametitle{Negasi Implikasi}
\begin{block}{Ekuivalensi}
\begin{align*}
p \rightarrow q &\equiv \neg p \vee q \\
\neg(p \rightarrow q) &\equiv p \wedge \neg q
\end{align*}
\end{block}
Negasi implikasi = anteseden benar DAN konsekuen salah.
\end{frame}

% ============================================================
% SECTION 3-6: Hirarki Prioritas Operator
% ============================================================
\section{Hirarki Prioritas Operator}

\begin{frame}
\frametitle{Urutan Prioritas (Tertinggi ke Terendah)}
\begin{table}
\small
\begin{tabular}{|c|l|l|}
\hline
\textbf{Prioritas} & \textbf{Simbol} & \textbf{Nama} \\ \hline
1 (tertinggi) & $\neg$ & Negasi \\ \hline
2 & $\wedge$ & Konjungsi \\ \hline
3 & $\vee$ & Disjungsi \\ \hline
4 & $\rightarrow$ & Implikasi \\ \hline
5 (terendah) & $\leftrightarrow$ & Biimplikasi \\ \hline
\end{tabular}
\end{table}
\end{frame}

\begin{frame}
\frametitle{Contoh Penguraian Ekspresi}
\begin{block}{Ekspresi: $\neg p \vee q \wedge r \rightarrow s \leftrightarrow t$}
\begin{enumerate}
  \item Negasi: $(\neg p)$
  \item Konjungsi: $(q \wedge r)$
  \item Disjungsi: $(\neg p) \vee (q \wedge r)$
  \item Implikasi: $\big((\neg p) \vee (q \wedge r)\big) \rightarrow s$
  \item Biimplikasi: $\Big(\big((\neg p) \vee (q \wedge r)\big) \rightarrow s\Big) \leftrightarrow t$
\end{enumerate}
\end{block}
\end{frame}

% ============================================================
% SECTION 3-7: Translasi Bahasa Alami
% ============================================================
\section{Translasi Bahasa Alami ke Logika Simbolik}

\begin{frame}
\frametitle{Konjungsi dalam Bahasa Alami}
\begin{contoh}
``Ibu memerintahkan adik menanak nasi, \textbf{tetapi} adik memilih menumis ayam.''\\
Kata ``tetapi'' mengindikasikan dua peristiwa sekaligus $\Rightarrow$ konjungsi.\\
$P \wedge T$ (dengan $P$ dan $T$ didefinisikan konsisten).
\end{contoh}
\end{frame}

\begin{frame}
\frametitle{``Jika'' dan ``Hanya Jika''}
\begin{block}{Aturan Praktis}
\begin{itemize}
  \item ``$A$ jika $B$'' $\Rightarrow$ $B \rightarrow A$
  \item ``$A$ hanya jika $B$'' $\Rightarrow$ $A \rightarrow B$
\end{itemize}
\end{block}
\begin{contoh}
``Anda dapat klaim garansi \textbf{hanya jika} kerusakan tidak menembus sirkuit.''\\
$K \rightarrow \neg R$
\end{contoh}
\end{frame}

% ============================================================
% RANGKUMAN
% ============================================================
\section{Rangkuman}

\begin{frame}
\frametitle{Rangkuman}
\begin{itemize}
  \item Proposisi: kalimat deklaratif dengan satu nilai kebenaran (T atau F).
  \item Kalimat terbuka mengandung variabel; non-proposisi: tanya, perintah, seru.
  \item Epistemologi: empiris vs derivasi deduktif.
  \item Operator: $\neg$, $\wedge$, $\vee$, $\rightarrow$, $\leftrightarrow$.
  \item Kontraposisi ekuivalen dengan implikasi asli; Konvers dan Invers tidak.
  \item Hukum De Morgan dan $\neg(p \rightarrow q) \equiv p \wedge \neg q$.
  \item Hirarki: $\neg$ $>$ $\wedge$ $>$ $\vee$ $>$ $\rightarrow$ $>$ $\leftrightarrow$.
\end{itemize}
\end{frame}

\begin{frame}
\frametitle{Kata Kunci}
\begin{center}
Proposisi $\bullet$ Kalimat Terbuka $\bullet$ Operator Logika $\bullet$ Implikasi $\bullet$ Biimplikasi\\
Kontraposisi $\bullet$ Hukum De Morgan $\bullet$ Hirarki Operator $\bullet$ Translasi Simbolik
\end{center}
\end{frame}

\end{document}
